% !PRINTED 67
\weberpage{67}
If, according to 5., one determines a function $T(u)$ with the zero points $\alpha_1, \alpha_2, \ldots, \alpha_m$, then:
\[
T(u)\varphi(u) = T_1(u)
\]
is a $T$-function of order $m$, and
\[
\varphi(u) = \frac{T_1(u)}{T(u)} \qquad \text{q.\,e.\,d.}
\]
The functions $T(u)$, $T_1(u)$ are \gesperrt{related}, since they have the same periodicity factors and therefore also the same characteristic.

The zero points of $T_1(u)$ are at the same time the zero points of $\varphi(u)$, and the sum of these \gesperrt{zero values} is therefore congruent with \gesperrt{the sum of the values of discontinuity}.

\begin{center}
\textbf{\S\,19. $T$-Functions of the First Order.}
\end{center}

We now proceed to determine more closely the $T$-functions of the first order, which we denote by $t(u)$, and from which, as we have seen, the remaining $T$-functions can be derived.

From 3. of the preceding section it follows that two $t$-functions of the same characteristic differ from one another only by a factor of the form
\[
Ce^{-\pi i(\lambda u^2+\mu u)}
\]
and we now wish, by an extension of the definition, to determine this factor still more closely.

This is to be done — which by § 17, (9) is possible without restriction of generality — in such a way that in the periodicity factors
\[
a_1 = 0, \quad b_1 = g_1, \quad b_2 = g_2
\]
holds, where $(g_1, g_2)$ is the characteristic; on account of the relation
\[
a_2\omega_1 - a_1\omega_2 = 1
\]
we then have
\[
a_2 = \frac{1}{\omega_1},
\]
and the conditions for this $t$-function then read (§ 17, 1.):
\weq{1}{\begin{aligned}
t(u+\omega_1) &= e^{-\pi i g_1}t(u),\\
t(u+\omega_2) &= e^{-\pi i \frac{2u+\omega_2}{\omega_1}}e^{-\pi i g_2}t(u),
\end{aligned}}
and by this the function $t(u)$ is determined up to a factor independent of $u$. In order to determine this factor still more closely, we
%CONT>>
